\section{Introducción}
La minería ilegal representa un problema ambiental, social y económico que no se limita a la extracción de oro, sino que incluye minerales como cobre, litio y cobalto, cuya demanda internacional ha impulsado actividades no autorizadas en diferentes regiones~\cite{boafo2025sustainability}. Estas actividades generan contaminación de fuentes hídricas, pérdida de bosques y afectaciones sobre comunidades locales e indígenas~\cite{mulenga2024aquatic,dplf2024illegal}. En la Amazonía, el aumento de la minería ilegal se ha relacionado con una mayor deforestación en territorios indígenas y con un incremento en la incidencia de malaria~\cite{santos2026illegal,amaral2024malaria}.

En Colombia, para 2022 se identificaron 94.733 hectáreas con evidencia de explotación de oro de aluvión, de las cuales el 73\,\% correspondía a explotación ilícita~\cite{minminas2023explotacion}. La situación es particularmente relevante en la Amazonía colombiana, donde las dificultades de acceso y la baja presencia institucional han favorecido su expansión~\cite{romero2020mineria,torres2020reimagining}. Las imágenes satelitales amplían la cobertura del monitoreo, pero la nubosidad frecuente dificulta la identificación de estas actividades~\cite{moomen2024assessing,becerra2024creating}. Las cámaras FLIR aerotransportadas complementan esa observación con videos de mayor detalle sobre zonas específicas, aunque producen grandes volúmenes de datos cuya revisión manual es lenta~\cite{delahoz2025enhancing}.

Como antecedente directo, la Fuerza Aeroespacial Colombiana (FAC) proporcionó videos obtenidos mediante sensores aerotransportados, y Acosta-Bernal et al.~\cite{acosta2025eyes} propusieron un pipeline basado en YOLOv11 para detectar indicadores asociados con minería y degradación ambiental. Ese trabajo dejó abiertas tres limitaciones que motivan este proyecto:

\begin{enumerate}
    \item \textbf{Calidad de la entrada.} Los videos incluyen simbología de vuelo superpuesta (HUD) que cubre parte de la escena, y ruido cuya forma cambia con el sensor y los filtros activos. Los trabajos previos usaron el preprocesamiento como apoyo de la detección, sin evaluar de manera independiente la eliminación del HUD ni la reducción del ruido~\cite{acosta2025eyes,delahoz2025enhancing}.
    \item \textbf{Fuga de datos.} Los fotogramas cercanos en el tiempo comparten escenario, objetos y perspectiva. Si se reparten al azar entre entrenamiento, validación y prueba, las métricas sobreestiman la capacidad de generalización~\cite{figueiredo2024leakage}.
    \item \textbf{Representación de la escena.} Las cajas delimitadoras localizan un elemento, pero no representan su extensión ni su relación con otros. En la línea base, las zonas de alteración del suelo de interés (SDZI) obtuvieron un recall de 0,425 y un mAP@50 de 0,402, muy por debajo del desempeño global~\cite{acosta2025eyes}.
\end{enumerate}

Este trabajo conjunto articula una solución para cada una: la limpieza de los videos (Laura Martínez-Galindo), el agrupamiento y particionamiento sin fuga de datos (Jorge Solórzano) y la segmentación panóptica con supervisión puntual (Manuela Galarza). Las tres etapas comparten los mismos datos y se encadenan: la limpieza entrega tres versiones del mismo conjunto de imágenes con nombres y etiquetas idénticos, el particionamiento define una única partición aplicable a las tres, y la segmentación panóptica se entrena y evalúa sobre ellas.

El resto del documento presenta el estado del arte, los objetivos, la metodología común, el desarrollo y los resultados de cada componente, la integración del pipeline y las conclusiones.
